% 2010 年陕西卷理科第 6 题：求样本平均数的程序框图。
% 来源 PDF：普通高考/2010/2010陕西理.pdf，SHA-256 8406e303e7aa1ac7813893781c903a7d9c904530028e6fd1c06062d1c32bd645。
% 原页：第 1 页（A3 三栏）第 1 栏下部，书页码 623。
% 重绘依据：照原页——竖排主干为 开始→输入 x1,…,x10→n=1,S=0→〔空白框〕→判断 n≥10?
%   →「是」向下进 x̄=S/n→输出 x̄→结束；「否」向左进 n=n+1 框，该框顶部再向上折回主干，
%   在空白框正上方汇入主干（原页此处只画一个向下箭头进入空白框）。空白框在原页为空。
\begin{tikzpicture}[semithick,
  every node/.style={inner sep=2pt},
  box/.style={draw, minimum width=3.5cm, minimum height=0.66cm, align=center},
  round/.style={draw, rounded corners=7pt, minimum width=1.7cm,
    minimum height=0.66cm, align=center},
  para/.style={draw, trapezium, trapezium left angle=72, trapezium right angle=108,
    minimum width=6.1cm, minimum height=0.74cm, align=center},
  dec/.style={draw, diamond, aspect=2.7, minimum width=2.9cm,
    minimum height=1.05cm, align=center},
  arr/.style={->, >=stealth}]

  \node[round] (start) at (0,0) {开始};
  \node[para]  (in)    at (0,-1.20) {输入 $x_1,\,x_2,\,\cdots,\,x_{10}$};
  \node[box]   (init)  at (0,-2.45) {$n=1$，$S=0$};
  \node[box, minimum height=0.86cm] (acc) at (0,-3.80) {};
  \node[dec]   (test)  at (0,-5.60) {$n\geqslant 10$};
  \node[box]   (inc)   at (-4.55,-3.80) {$n=n+1$};
  \node[box, minimum height=1.50cm, text height=0.72cm, text depth=0.62cm] (avg) at (0,-7.45) {$\bar x=\dfrac{S}{n}$};
  \node[para]  (out)   at (0,-9.55) {输出 $\bar x$};
  \node[round] (end)   at (0,-11.20) {结束};

  \draw[arr] (start) -- (in);
  \draw[arr] (in) -- (init);
  \draw[arr] (init) -- (acc);

  \draw[arr] (acc) -- (test);
  \draw[arr] (test) -- (avg);
  \node[right] at (0.16,-6.40) {是};
  \draw[arr] (avg) -- (out);
  \draw[arr] (out) -- (end);

  % 「否」分支：由判断框左顶点向左，再向上进入 n=n+1 框
  \draw[arr] (test.west) -- (-4.55,-5.60) -- (inc.south);
  \node[above] at (-2.55,-5.45) {否};

  % 回边：n=n+1 框顶部向上、再向右，汇入空白框上方的竖线
  \draw (inc.north) -- (-4.55,-2.95) -- (0,-2.95);
\end{tikzpicture}
